\documentclass[11pt]{article}

%==============================================================================
%  A blueprint for the formalisation of the moduli stack of elliptic curves
%  via fibered categories in Lean 4 / Mathlib.
%
%  The document is written so that it compiles standalone (the leanblueprint
%  macros \lean, \leanok, \uses, \notready, \discussion are defined below as
%  no-ops or lightweight markers).  To use it inside a genuine `leanblueprint'
%  project, delete the "STANDALONE SHIMS" block and let plastexdepgraph supply
%  the real definitions.
%==============================================================================

\usepackage[margin=1in]{geometry}
\usepackage{amsmath,amssymb,amsthm,mathtools}
\usepackage{stmaryrd}
\usepackage{enumitem}
\usepackage{hyperref}
\usepackage{xcolor}
\hypersetup{colorlinks=true,linkcolor=blue!55!black,citecolor=green!45!black,
            urlcolor=blue!55!black}

%---------------------------- STANDALONE SHIMS --------------------------------
\newcommand{\lean}[1]{\par\smallskip\noindent{\footnotesize\texttt{Lean:\ #1}}\par}
\newcommand{\leanok}{{\footnotesize\textcolor{green!50!black}{\ [in Mathlib / formalised]}}}
\newcommand{\uses}[1]{}
\newcommand{\notready}{{\footnotesize\textcolor{red!60!black}{\ [not ready]}}}
\newcommand{\discussion}[1]{}
%------------------------------------------------------------------------------

\theoremstyle{definition}
\newtheorem{definition}{Definition}[section]
\newtheorem{construction}[definition]{Construction}
\newtheorem{convention}[definition]{Convention}
\newtheorem{notation}[definition]{Notation}
\theoremstyle{plain}
\newtheorem{lemma}[definition]{Lemma}
\newtheorem{proposition}[definition]{Proposition}
\newtheorem{theorem}[definition]{Theorem}
\newtheorem{corollary}[definition]{Corollary}
\theoremstyle{remark}
\newtheorem{remark}[definition]{Remark}
\newtheorem{warning}[definition]{Warning}

\newcommand{\Aff}{\mathrm{Aff}}
\newcommand{\CRing}{\mathrm{CommRing}}
\newcommand{\CRingCat}{\mathbf{CommRingCat}}
\newcommand{\Sch}{\mathrm{Sch}}
\newcommand{\Cat}{\mathbf{Cat}}
\newcommand{\Grpd}{\mathbf{Grpd}}
\newcommand{\Spec}{\operatorname{Spec}}
\newcommand{\op}{\mathrm{op}}
\newcommand{\Weier}{\mathcal{W}}
\newcommand{\Ell}{\mathcal{E}\!\ell\!\ell}
\newcommand{\MW}{\mathcal{M}^{\mathrm{W}}_{1,1}}
\newcommand{\Mss}{\mathcal{M}_{1,1}}
\newcommand{\VC}{\mathrm{VC}}
\newcommand{\Isom}{\underline{\mathrm{Isom}}}
\newcommand{\Gm}{\mathbb{G}_{\mathrm m}}
\newcommand{\Ga}{\mathbb{G}_{\mathrm a}}
\newcommand{\id}{\mathrm{id}}
\newcommand{\Hom}{\operatorname{Hom}}
\newcommand{\map}{\mathrm{map}}
\newcommand{\isEll}{\mathtt{IsElliptic}}
\newcommand{\isHomLift}{\mathtt{IsHomLift}}
\newcommand{\isCart}{\mathtt{IsCartesian}}
\newcommand{\isSCart}{\mathtt{IsStronglyCartesian}}
\newcommand{\isFib}{\mathtt{IsFibered}}
\newcommand{\onel}{\mathbf 1}
\renewcommand{\hom}{\longrightarrow}

\title{\bfseries A blueprint for the moduli stack of elliptic curves\\
       via fibered categories in Mathlib}
\author{}
\date{}

\begin{document}
\maketitle

\begin{abstract}
We give a detailed blueprint for the formalisation, in Lean~4/Mathlib, of the
moduli stack of elliptic curves through the language of fibered categories.
The construction proceeds in five parts.  Part~\ref{sec:substrate} records the
Mathlib substrate.  Part~\ref{sec:base} fixes the base site (affine schemes with
the fppf topology).  Part~\ref{sec:weier} builds the fibered category
$\Weier\to\Aff$ of Weierstrass curves out of the existing
\texttt{WeierstrassCurve} and \texttt{WeierstrassCurve.VariableChange} API, and
proves that its projection is fibered in groupoids.  Part~\ref{sec:moduli}
extracts the full subcategory $\MW$ of elliptic objects and records the quotient
presentation $\MW\simeq[\Spec A/G]$ with
$A=\mathbb Z[a_1,a_2,a_3,a_4,a_6][\Delta^{-1}]$ and $G=\Ga^3\rtimes\Gm$.
Part~\ref{sec:descent} isolates the descent (stack) condition and reduces it to
faithfully flat descent statements, several of which are not yet in Mathlib.
Part~\ref{sec:comparison} compares $\MW$ with the geometric stack $\Mss$ of
proper smooth genus-one curves with a section and states the comparison theorem,
whose hypotheses expose the scheme-theoretic prerequisites still missing from
Mathlib.  A milestone ladder (Part~\ref{sec:ladder}) orders the work by
Mathlib-readiness.

\medskip\noindent\emph{Correctness note.}  Throughout, declarations already
present in Mathlib are marked \leanok\ and annotated with their module path; all
other items are targets.  Every nontrivial target carries an explicit list of
its prerequisites via \texttt{\textbackslash uses}, so that the dependency graph
is complete.
\end{abstract}

\tableofcontents

%==============================================================================
\section{Conventions and the Mathlib substrate}\label{sec:substrate}
%==============================================================================

\begin{convention}\label{conv:univ}
We work with Grothendieck universes as Mathlib does.  A \emph{category} means a
term of \texttt{CategoryTheory.Category.\{v,u\}}; morphisms are written
$X\hom Y$, composition $f\gg g$ (``$f$ then $g$''), identities $\onel_X$.
Opposite categories are $\mathcal C^{\op}$.  For a functor $p:\mathcal X\to\mathcal S$
we call $\mathcal S$ the \emph{base} and $\mathcal X$ the \emph{total category}.
\end{convention}

\begin{convention}[Rings and affine schemes]\label{conv:aff}
$\CRingCat$ is the category of commutative rings
(\texttt{Mathlib.Algebra.Category.Ring.Basic}).  We set
$\Aff\coloneqq\CRingCat^{\op}$ and identify it with the category of affine
schemes; an object of $\Aff$ is a ring $R$ (written $\Spec R$ when the geometric
reading is intended), and a morphism $\Spec R\to\Spec R'$ in $\Aff$ is a ring
homomorphism $f\colon R'\to R$, denoted $\underline f$ when regarded as an arrow
of $\Aff$.  \emph{Base change} along $\underline f$ is application of $f$.
\lean{CommRingCat}\leanok
\end{convention}

\subsection{The fibered-category API in Mathlib}

The following are present in \texttt{Mathlib.CategoryTheory.FiberedCategory.*};
we recall the interface we will use.

\begin{definition}[Lift of a morphism]\label{def:homlift}
For $p:\mathcal X\to\mathcal S$, objects $a,b\in\mathcal X$ over $R,S\in\mathcal S$
(i.e.\ $p(a)=R$, $p(b)=S$), a base arrow $f:R\to S$ and $\varphi:a\hom b$, the
proposition $p.\isHomLift\,f\,\varphi$ asserts that $p(\varphi)=f$ modulo the
equalities $p(a)=R$, $p(b)=S$ (implemented with \texttt{eqToHom}).
\lean{CategoryTheory.Functor.IsHomLift}\leanok
\end{definition}

\begin{definition}[Cartesian and strongly cartesian morphisms]\label{def:cartesian}
With notation as in Definition~\ref{def:homlift} and $\varphi$ a lift of
$f:R\to S$:
\begin{enumerate}[label=(\alph*)]
\item $\varphi$ is \emph{cartesian} ($p.\isCart\,f\,\varphi$) if for every
$a'\in\mathcal X$ over $R$ and every $\psi:a'\hom b$ lifting $f$ there is a
unique vertical $\chi:a'\hom a$ (a lift of $\onel_R$) with $\chi\gg\varphi=\psi$.
\item $\varphi$ is \emph{strongly cartesian} ($p.\isSCart\,f\,\varphi$) if for
every $a'\in\mathcal X$, every $g:p(a')\to R$ and every $\psi:a'\hom b$ lifting
$f\circ g$ there is a unique $\chi:a'\hom a$ lifting $g$ with
$\chi\gg\varphi=\psi$.
\end{enumerate}
\lean{CategoryTheory.Functor.IsCartesian, CategoryTheory.Functor.IsStronglyCartesian}\leanok
\end{definition}

\begin{lemma}[Basic API for (strongly) cartesian morphisms]\label{lem:scart-api}
Strongly cartesian morphisms are stable under composition; a strongly cartesian
morphism lifting an isomorphism is an isomorphism; two strongly cartesian lifts
of the same arrow with the same target are related by a unique vertical
isomorphism of their sources.  All three are available in Mathlib.
\lean{CategoryTheory.Functor.IsStronglyCartesian.comp,
      CategoryTheory.Functor.IsStronglyCartesian.isIso\_of\_base\_isIso,
      CategoryTheory.Functor.IsStronglyCartesian.domainUniqueUpToIso}\leanok
\uses{def:cartesian}
\end{lemma}

\begin{definition}[Fibered functor]\label{def:isfibered}
$p:\mathcal X\to\mathcal S$ is \emph{prefibered} ($p.\mathtt{IsPreFibered}$) if
for every $f:R\to S$ in $\mathcal S$ and every $b$ over $S$ there is a strongly
cartesian $\varphi:a\hom b$ lifting $f$; it is \emph{fibered}
($p.\isFib$) if in addition strongly cartesian morphisms are stable under
composition (automatic from Lemma~\ref{lem:scart-api}).  Mathlib packages the
prefibered pullback as \texttt{pullbackObj} and \texttt{pullbackMap}.
\lean{CategoryTheory.Functor.IsPreFibered, CategoryTheory.Functor.IsFibered,
      CategoryTheory.Functor.pullbackObj, CategoryTheory.Functor.pullbackMap}\leanok
\uses{def:cartesian,lem:scart-api}
\end{definition}

\begin{definition}[Fibre categories]\label{def:fiber}
For $S\in\mathcal S$, $p.\mathtt{Fiber}\,S$ is the subtype
$\{a:\mathcal X\mid p(a)=S\}$, made a category with morphisms the vertical
morphisms (lifts of $\onel_S$).  The inclusion $\mathtt{fiberInclusion}$ into
$\mathcal X$ is faithful, and every $\varphi$ in the fibre satisfies
$p.\isHomLift\,\onel_S\,\varphi$.
\lean{CategoryTheory.Functor.Fiber, CategoryTheory.Functor.Fiber.fiberInclusion}\leanok
\uses{def:homlift}
\end{definition}

\begin{definition}[Fibered in groupoids]\label{def:cfg}
A fibered functor $p$ is \emph{fibered in groupoids} (a \emph{CFG}) when every
morphism of $\mathcal X$ is strongly cartesian; equivalently every fibre
$p.\mathtt{Fiber}\,S$ is a groupoid and $p$ is prefibered.  The equivalence of
the two formulations is standard \cite[Def.~3.21, \S3.6]{FGA} and is the target
lemma \texttt{isFibered\_iff\_fibers\_groupoid} below (see
Proposition~\ref{prop:cfg-criterion}).
\uses{def:isfibered,def:fiber}
\end{definition}

\begin{remark}[The Grothendieck and CoGrothendieck constructions]\label{rem:grothendieck}
Mathlib (\texttt{Mathlib.CategoryTheory.Bicategory.Grothendieck}) provides two
Grothendieck constructions for a pseudofunctor valued in $\Cat$
\cite{Vistoli2008}.  For a covariant
$F:\mathtt{LocallyDiscrete}\,\mathcal S\to\Cat$ it provides $\int F$
(\texttt{Pseudofunctor.Grothendieck}), whose objects are pairs $(S,a)$ with
$a\in F(S)$ and whose morphisms $(R,b)\to(S,a)$ are pairs $(f,h)$ with
$f:R\to S$ in $\mathcal S$ and $h:F(f)(a)\to b$.  For a contravariant
$F:\mathtt{LocallyDiscrete}\,\mathcal S^{\op}\to\Cat$ it provides
$\int^{\mathrm c}\!F$ (\texttt{Pseudofunctor.CoGrothendieck}), whose objects are
pairs $(S,a)$ with $a\in F(S)$ and whose morphisms $(R,b)\to(S,a)$ are pairs
$(f,h)$ with $f:R\to S$ in $\mathcal S$ and $h:b\to F(f)(a)$, together with the
projection $\mathtt{forget}\,F:\int^{\mathrm c}\!F\to\mathcal S$,
$(S,a)\mapsto S$, $(f,h)\mapsto f$.  The Mathlib naming convention reserves
\emph{Grothendieck} for the covariant construction and \emph{CoGrothendieck} for
the contravariant one, consistent with \texttt{CategoryTheory.Grothendieck} for
$1$-functors.  The moduli problem is contravariant in $\mathcal S=\Aff$, so the
relevant construction is $\int^{\mathrm c}\!F$; see
Construction~\ref{con:pseudofunctor}.
\lean{CategoryTheory.Pseudofunctor.Grothendieck, CategoryTheory.Pseudofunctor.CoGrothendieck}\leanok
\end{remark}

\begin{proposition}[The CoGrothendieck construction is a fibered category]\label{prop:cogroth-fibered}
For $F:\mathtt{LocallyDiscrete}\,\mathcal S^{\op}\to\Cat$ the projection
$\mathtt{forget}\,F:\int^{\mathrm c}\!F\to\mathcal S$ is a fibered category:
\begin{enumerate}[label=(\alph*)]
\item for $f:R\to S$ in $\mathcal S$ and $a\in F(S)$, the arrow
$\mathtt{cartesianLift}\,a\,f$ with domain $(R,\,F(f)(a))$ is strongly cartesian;
\item $\mathtt{forget}\,F$ carries an \texttt{IsFibered} instance;
\item $\mathtt{forget}\,F$ carries a \texttt{HasFibers} instance whose fibre over
$S$ is $F(S)$, and the induced functor
$F(S)\to(\mathtt{forget}\,F).\mathtt{Fiber}\,S$ is an equivalence.
\end{enumerate}
All three are established in
\texttt{Mathlib.CategoryTheory.FiberedCategory.Grothendieck} (using
\texttt{FiberedCategory.HasFibers}).
\lean{CategoryTheory.Pseudofunctor.CoGrothendieck.isStronglyCartesian\_homCartesianLift,
      CategoryTheory.Pseudofunctor.CoGrothendieck.instIsFiberedForget,
      CategoryTheory.Pseudofunctor.CoGrothendieck.instHasFibersForget}\leanok
\uses{rem:grothendieck,def:isfibered,def:fiber}
\end{proposition}

\subsection{Sites and sheaves in Mathlib}

\begin{definition}[Grothendieck topology, sheaf]\label{def:site}
Mathlib has sieves, Grothendieck topologies (maximality, stability under
pullback, transitivity), pretopologies, coverages, the sheaf condition for
presheaves valued in an arbitrary category, sheafification for suitable targets,
and the comparison lemma between pretopologies and topologies.
\lean{CategoryTheory.GrothendieckTopology, CategoryTheory.Pretopology,
      CategoryTheory.Sieve, CategoryTheory.Presheaf.IsSheaf,
      CategoryTheory.Sheaf}\leanok
\end{definition}

\begin{warning}[Stacks are not in Mathlib]\label{warn:no-stacks}
As of the reference commit, Mathlib contains \emph{sheaf} theory (a $1$-truncated
descent theory) but no notion of \emph{stack} (descent for a fibered category),
no stackification, and no fppf/étale topology on schemes as a named object.  All
of Part~\ref{sec:descent} that concerns descent \emph{of objects} is therefore
new content.  This is the principal gap between the present blueprint and a
finished formalisation.
\end{warning}

\subsection{The elliptic-curve API in Mathlib}\label{ss:ecapi}

\begin{definition}[Weierstrass curve]\label{def:weier}
For $R:\CRing$, $\texttt{WeierstrassCurve}\,R$ is the five-tuple
$(a_1,a_2,a_3,a_4,a_6)\in R^5$.  Derived quantities $b_2,b_4,b_6,b_8$,
$c_4,c_6$, and the discriminant $\Delta$ are defined, with
$c_4^3-c_6^2=1728\,\Delta$ and $\Delta=(-b_2^2b_8-8b_4^3-27b_6^2+9b_2b_4b_6)$.
Base change along $f\colon R\to A$ is $W.\map\,f:\texttt{WeierstrassCurve}\,A$,
obtained by applying $f$ to each coefficient.
\lean{WeierstrassCurve, WeierstrassCurve.Delta, WeierstrassCurve.map}\leanok
\end{definition}

\begin{lemma}[Functoriality of base change]\label{lem:map-func}
$W.\map\,(\mathrm{RingHom.id}\,R)=W$ and
$(W.\map\,f).\map\,g=W.\map\,(g\circ f)$; moreover
$(W.\map\,f).\Delta=f(W.\Delta)$ and likewise for $b_i,c_4,c_6$.  Injectivity of
$W\mapsto W.\map\,f$ for injective $f$ is \texttt{map\_injective}.
\lean{WeierstrassCurve.map\_id, WeierstrassCurve.map\_map,
      WeierstrassCurve.map\_injective}\leanok
\uses{def:weier}
\end{lemma}

\begin{definition}[Variable changes]\label{def:vc}
$\texttt{WeierstrassCurve.VariableChange}\,R$ is the datum
$C=(u,r,s,t)$ with $u\in R^{\times}$ and $r,s,t\in R$.  These form a group under
the classical composition law \cite[III, Table~3.1]{Silverman}, and act on
$\texttt{WeierstrassCurve}\,R$ by
\[
  C\cdot W=\bigl(a_1',a_2',a_3',a_4',a_6'\bigr),
\]
where, writing $u^{-1}\in R^{\times}$,
\[
\begin{aligned}
 a_1'&=u^{-1}(a_1+2s), &
 a_2'&=u^{-2}(a_2-sa_1+3r-s^2),\\
 a_3'&=u^{-3}(a_3+ra_1+2t), &
 a_4'&=u^{-4}(a_4-sa_3+2ra_2-(t+rs)a_1+3r^2-2st),\\
 a_6'&=u^{-6}(a_6+ra_4+r^2a_2+r^3-ta_3-t^2-rta_1). &&
\end{aligned}
\]
This is a left action ($\onel\cdot W=W$, $(C_1C_2)\cdot W=C_1\cdot(C_2\cdot W)$).
\lean{WeierstrassCurve.VariableChange, WeierstrassCurve.instGroupVariableChange,
      WeierstrassCurve.instMulActionVariableChange}\leanok
\end{definition}

\begin{lemma}[Transformation of invariants]\label{lem:vc-invariants}
For $C=(u,r,s,t)$ one has
$(C\cdot W).\Delta=u^{-12}\,W.\Delta$, $(C\cdot W).c_4=u^{-4}\,W.c_4$,
$(C\cdot W).c_6=u^{-6}\,W.c_6$.
\lean{WeierstrassCurve.variableChange\_Delta, WeierstrassCurve.variableChange\_c4,
      WeierstrassCurve.variableChange\_c6}\leanok
\uses{def:vc,def:weier}
\end{lemma}

\begin{definition}[Elliptic curve; $j$-invariant]\label{def:iselliptic}
$\texttt{WeierstrassCurve.IsElliptic}\,W$ is the class asserting that $\Delta$ is
a unit; under it $\Delta'\in R^{\times}$ is the unit with
$(\Delta':R)=\Delta$, and the $j$-invariant is $j\coloneqq (\Delta')^{-1}c_4^3$.
Every $j_0\in F$ over a field $F$ is realised by an explicit curve
$\texttt{ofJ}\,j_0$.
\lean{WeierstrassCurve.IsElliptic, WeierstrassCurve.Delta', WeierstrassCurve.coe\_Delta',
      WeierstrassCurve.j, WeierstrassCurve.ofJ}\leanok
\uses{def:weier}
\end{definition}

\begin{warning}[The affine model is a presentation, not $\Mss$]\label{warn:pic}
The Mathlib module docstring records the decisive caveat: the type
$\texttt{WeierstrassCurve}\,R$ with $\isEll$ ``gives a type which can be beefed up
to a category which is equivalent to the category of elliptic curves over
$\Spec R$ in the case that $R$ has trivial Picard group $\mathrm{Pic}(R)$ or,
slightly more generally, when its $12$-torsion is trivial.''  In general the two
differ by line-bundle twists; see Part~\ref{sec:comparison}.  The object built in
Parts~\ref{sec:weier}--\ref{sec:moduli} is therefore the \emph{Weierstrass
presentation} of $\Mss$, whose fppf-stackification is $\Mss$.
\end{warning}

%==============================================================================
\section{The base site: affine schemes with the fppf topology}\label{sec:base}
%==============================================================================

\begin{definition}[fppf covers of an affine scheme]\label{def:fppf-cover}
A \emph{standard fppf cover} of $R\in\Aff$ is a finite family of ring maps
$\{R\to R_i\}_{i\in I}$, each flat and of finite presentation, such that
$R\to\prod_i R_i$ is faithfully flat (equivalently $\Spec(\prod R_i)\to\Spec R$
is surjective).  The associated sieve on $\Spec R$ in $\Aff$ is generated by the
$\underline{(R\to R_i)}:\Spec R_i\to\Spec R$.
\lean{}\notready
\end{definition}

\begin{lemma}[The standard fppf data form a pretopology]\label{lem:fppf-pretop}
The families of Definition~\ref{def:fppf-cover} contain the isomorphisms, are
stable under base change (pullback of a flat finite-presentation map is again
such, and faithful flatness is preserved), and are stable under composition
(refinement).  Hence they define a \texttt{Pretopology} on $\Aff$, whose
associated Grothendieck topology is the \emph{fppf topology} $\tau_{\mathrm{fppf}}$.
\lean{}\notready
\uses{def:fppf-cover,def:site}
\end{lemma}
\begin{proof}[Proof sketch]
Pullback in $\Aff$ is the tensor product: the base change of $R\to R_i$ along
$R\to R'$ is $R'\to R'\otimes_R R_i$.  Flatness and finite presentation are
stable under base change (\texttt{Module.Flat.baseChange},
\texttt{Algebra.FinitePresentation}\,base-change lemmas).  Faithful flatness of
$R\to\prod R_i$ pulls back to faithful flatness of $R'\to\prod R'\otimes_R R_i$
(\texttt{Module.FaithfullyFlat}\,base-change).  Stability under composition uses
that a composite of flat (resp.\ finitely presented, resp.\ faithfully flat) maps
is again such.  The pretopology axioms of
\texttt{CategoryTheory.Pretopology} then hold by construction.  The Zariski and
étale (co)topologies arise by restricting to open immersions
(localisations $R\to R_f$) resp.\ to étale maps; they are coarser than
$\tau_{\mathrm{fppf}}$.
\end{proof}

\begin{remark}[Set-theoretic size]\label{rem:small-site}
To stay within Mathlib's universe discipline one works either with the
\emph{small} fppf site of a fixed base ring $R_0$
(the slice $\Aff_{/\Spec R_0}$ of finitely presented $R_0$-algebras) or with a
suitable small category of rings.  The moduli constructions below are
insensitive to this choice; we write $\Aff$ generically and fix a small model
when formalising.
\end{remark}

%==============================================================================
\section{The fibered category of Weierstrass curves}\label{sec:weier}
%==============================================================================

Throughout this section base arrows are read as in Convention~\ref{conv:aff}:
an $\Aff$-arrow $\underline f:\Spec R\to\Spec R'$ is a ring map $f:R'\to R$.

\subsection{Equivariance of the action under base change}

\begin{definition}[Base change of a variable change]\label{def:vc-map}
For a ring homomorphism $f:R\to A$ define $\map\,f:\VC(R)\to\VC(A)$ on underlying
data by $(u,r,s,t)\mapsto(\mathrm{Units.map}\,f\,u,\ f(r),\ f(s),\ f(t))$.
\lean{WeierstrassCurve.VariableChange.map}\notready
\uses{def:vc}
\end{definition}

\begin{lemma}[$\map\,f$ preserves the identity]\label{lem:vc-map-one}
For $f:R\to A$, $\map\,f\,\onel=\onel$ in $\VC(A)$.
\lean{WeierstrassCurve.VariableChange.map\_one}\notready
\uses{def:vc-map,def:vc}
\end{lemma}

\begin{lemma}[$\map\,f$ preserves the product]\label{lem:vc-map-mul}
For $f:R\to A$ and $C_1,C_2\in\VC(R)$,
$\map\,f\,(C_1C_2)=(\map\,f\,C_1)(\map\,f\,C_2)$; hence $\map\,f$ is a group
homomorphism $\VC(R)\to\VC(A)$.
\lean{WeierstrassCurve.VariableChange.map\_mul}\notready
\uses{def:vc-map,def:vc}
\end{lemma}

\begin{lemma}[Functoriality of $\map$: identity]\label{lem:vc-map-id}
$\map\,(\mathrm{RingHom.id}\,R)=\mathrm{id}_{\VC(R)}$.
\lean{WeierstrassCurve.VariableChange.map\_id}\notready
\uses{def:vc-map}
\end{lemma}

\begin{lemma}[Functoriality of $\map$: composition]\label{lem:vc-map-comp}
For $f:R\to A$ and $g:A\to B$ one has $\map\,(g\circ f)=\map\,g\circ\map\,f$;
equivalently $\map\,g\,(\map\,f\,C)=\map\,(g\circ f)\,C$ for all $C\in\VC(R)$.
\lean{WeierstrassCurve.VariableChange.map\_map}\notready
\uses{def:vc-map}
\end{lemma}

\begin{lemma}[Equivariance of the action under base change]\label{lem:equivariance}
For $f:R\to A$, $C\in\VC(R)$ and $W:\texttt{WeierstrassCurve}\,R$,
\[
   (C\cdot W).\map\,f \;=\; (\map\,f\,C)\cdot(W.\map\,f).
\]
\lean{WeierstrassCurve.map\_variableChange}\notready
\uses{def:vc-map,def:weier,lem:map-func}
\end{lemma}
\begin{proof}[Proof sketch]
Both sides are Weierstrass curves over $A$; by \texttt{WeierstrassCurve.ext} it
suffices to compare the five coefficients $a_1,\dots,a_6$.  Each coefficient of
the left side is $f$ applied to the corresponding coefficient of $C\cdot W$, i.e.\
$f$ applied to a polynomial in $u^{-1},r,s,t,a_i$ (Definition~\ref{def:vc}); each
coefficient of the right side is the same polynomial evaluated at
$\mathrm{Units.map}\,f\,u^{-1}=f(u^{-1})$, $f(r),f(s),f(t),f(a_i)$.  The two agree
because $f$ is a ring homomorphism and $\mathrm{Units.map}\,f$ commutes with
inversion; each coefficient identity is closed by \texttt{map\_mul},
\texttt{map\_add}, \texttt{map\_pow}, \texttt{Units.coe\_map} and \texttt{ring}.
\end{proof}

\subsection{The total category}

\begin{definition}[Objects of $\Weier$]\label{def:total-obj}
An object of $\Weier$ is a pair $(R,W)$ with $R\in\Aff$ and
$W:\texttt{WeierstrassCurve}\,R$.
\lean{}\notready
\uses{def:weier,conv:aff}
\end{definition}

\begin{definition}[Morphisms of $\Weier$]\label{def:total-hom}
A morphism $(R,W)\to(R',W')$ is a pair $(\underline f,C)$ where
$\underline f:\Spec R\to\Spec R'$ is an $\Aff$-arrow (a ring map $f:R'\to R$) and
$C\in\VC(R)$ satisfies
\begin{equation}\label{eq:mor}
   C\cdot W \;=\; W'.\map\,f .
\end{equation}
Two morphisms are equal iff their $\Aff$-arrows and their $\VC$-components agree,
the condition \eqref{eq:mor} being a proposition.
\lean{}\notready
\uses{def:total-obj,def:vc,def:vc-map,conv:aff}
\end{definition}

\begin{lemma}[The identity datum is a morphism]\label{lem:id-welldef}
For $(R,W)\in\Weier$ the pair $(\underline{\id_R},\onel)$ satisfies \eqref{eq:mor}:
$\onel\cdot W=W=W.\map\,(\mathrm{RingHom.id}\,R)$.
\lean{}\notready
\uses{def:total-hom,lem:map-func,def:vc}
\end{lemma}

\begin{definition}[Identity morphism]\label{def:total-id}
The identity of $(R,W)$ is $\onel_{(R,W)}\coloneqq(\underline{\id_R},\onel)$,
well-defined by Lemma~\ref{lem:id-welldef}.
\lean{}\notready
\uses{lem:id-welldef}
\end{definition}

\begin{lemma}[The composite datum is a morphism]\label{lem:comp-welldef}
Given
$(R,W)\xrightarrow{(\underline f,C)}(R',W')\xrightarrow{(\underline g,C')}(R'',W'')$
(so $f:R'\to R$, $g:R''\to R'$), the pair
$(\underline{f\circ g},\ (\map\,f\,C')\cdot C)$ satisfies \eqref{eq:mor} over $R$:
\[
   \bigl((\map\,f\,C')\cdot C\bigr)\cdot W=W''.\map\,(f\circ g).
\]
\lean{}\notready
\uses{def:total-hom,lem:equivariance,lem:map-func}
\end{lemma}
\begin{proof}
Using $C\cdot W=W'.\map\,f$ and $C'\cdot W'=W''.\map\,g$,
\[
\begin{aligned}
 ((\map\,f\,C')\cdot C)\cdot W
   &=(\map\,f\,C')\cdot(C\cdot W)          &&\text{(\texttt{mul\_smul})}\\
   &=(\map\,f\,C')\cdot(W'.\map\,f)         &&\text{(hypothesis on }C)\\
   &=(C'\cdot W').\map\,f                    &&\text{(Lemma~\ref{lem:equivariance})}\\
   &=(W''.\map\,g).\map\,f                   &&\text{(hypothesis on }C')\\
   &=W''.\map\,(f\circ g)                    &&\text{(Lemma~\ref{lem:map-func}).}
\end{aligned}
\]
\end{proof}

\begin{definition}[Composition]\label{def:total-comp}
The composite of $(\underline f,C)$ and $(\underline g,C')$ as above is
\begin{equation}\label{eq:comp}
   (\underline f,C)\gg(\underline g,C')\;\coloneqq\;
   (\underline{f\circ g},\ (\map\,f\,C')\cdot C ),
\end{equation}
well-defined by Lemma~\ref{lem:comp-welldef}.
\lean{}\notready
\uses{lem:comp-welldef,def:vc-map}
\end{definition}

\begin{lemma}[Left identity]\label{lem:total-id-comp}
$\onel_{(R,W)}\gg(\underline f,C)=(\underline f,C)$ for every
$(\underline f,C):(R,W)\to(R',W')$.
\lean{}\notready
\uses{def:total-comp,def:total-id,lem:vc-map-id}
\end{lemma}
\begin{proof}
The $\Aff$-component is $\underline{\id_R\circ f}=\underline f$; the
$\VC$-component is $(\map\,\id\,C)\cdot\onel=C\cdot\onel=C$, using
Lemma~\ref{lem:vc-map-id} and \texttt{mul\_one}.
\end{proof}

\begin{lemma}[Right identity]\label{lem:total-comp-id}
$(\underline f,C)\gg\onel_{(R',W')}=(\underline f,C)$.
\lean{}\notready
\uses{def:total-comp,def:total-id,lem:vc-map-one}
\end{lemma}
\begin{proof}
The $\Aff$-component is $\underline{f\circ\id_{R'}}=\underline f$; the
$\VC$-component is $(\map\,f\,\onel)\cdot C=\onel\cdot C=C$, using
Lemma~\ref{lem:vc-map-one} and \texttt{one\_mul}.
\end{proof}

\begin{lemma}[Associativity]\label{lem:total-assoc}
For composable $(\underline f,C),(\underline g,C'),(\underline h,C'')$,
\[
 ((\underline f,C)\gg(\underline g,C'))\gg(\underline h,C'')
 =(\underline f,C)\gg((\underline g,C')\gg(\underline h,C'')).
\]
\lean{}\notready
\uses{def:total-comp,lem:vc-map-mul,lem:vc-map-comp}
\end{lemma}
\begin{proof}
Both $\Aff$-components equal $\underline{f\circ g\circ h}$.  The $\VC$-components
are
\[
 (\map\,(f\circ g)\,C'')\cdot((\map\,f\,C')\cdot C)
 \quad\text{and}\quad
 (\map\,f\,((\map\,g\,C'')\cdot C'))\cdot C .
\]
Expanding the second by Lemma~\ref{lem:vc-map-mul} ($\map\,f$ preserves products)
and Lemma~\ref{lem:vc-map-comp}
($\map\,f\,(\map\,g\,C'')=\map\,(f\circ g)\,C''$) yields
$((\map\,(f\circ g)\,C'')\cdot(\map\,f\,C'))\cdot C$, equal to the first by
\texttt{mul\_assoc}.
\end{proof}

\begin{definition}[The category $\Weier$]\label{def:total}
The objects (Definition~\ref{def:total-obj}), morphisms
(Definition~\ref{def:total-hom}), identities (Definition~\ref{def:total-id}) and
composition (Definition~\ref{def:total-comp}), together with
Lemmas~\ref{lem:total-id-comp}, \ref{lem:total-comp-id} and
\ref{lem:total-assoc}, constitute a category $\Weier$.
\lean{Mathlib.AlgebraicGeometry.EllipticCurve.Moduli.Weierstrass}\notready
\uses{def:total-obj,def:total-hom,def:total-id,def:total-comp,lem:total-id-comp,lem:total-comp-id,lem:total-assoc}
\end{definition}

\begin{definition}[The projection $p_{\Weier}$]\label{def:proj}
$p_{\Weier}:\Weier\to\Aff$ sends $(R,W)\mapsto\Spec R$ and
$(\underline f,C)\mapsto\underline f$.  It preserves identities and composition
by Definitions~\ref{def:total-id} and \ref{def:total-comp}, whose $\Aff$-components
are $\underline{\id}$ and the composite of the $\Aff$-components respectively.
\lean{}\notready
\uses{def:total}
\end{definition}

\subsection{Cartesian structure}

\begin{lemma}[Morphisms are homlifts]\label{lem:homlift-mor}
Every morphism $(\underline f,C):(R,W)\to(R',W')$ satisfies
$p_{\Weier}.\isHomLift\,\underline f\,(\underline f,C)$, by definition of the
projection.
\lean{}\notready
\uses{def:proj,def:homlift}
\end{lemma}

\begin{lemma}[Vertical morphisms]\label{lem:vertical-iff}
A morphism $(\underline g,C):(R,W)\to(R,W_0)$ is vertical (a lift of $\onel_R$)
iff $g=\id_R$; the vertical morphisms $(R,W)\to(R,W_0)$ are then exactly the
$C\in\VC(R)$ with $C\cdot W=W_0$.
\lean{}\notready
\uses{lem:homlift-mor,def:homlift,def:fiber}
\end{lemma}

\begin{lemma}[Fibre hom-sets]\label{lem:fiber-hom}
For $R\in\Aff$ and $W,W_0$ over $R$, the hom-set in the fibre
$p_{\Weier}.\mathtt{Fiber}\,(\Spec R)$ is
$\Hom(W,W_0)=\{\,C\in\VC(R)\mid C\cdot W=W_0\,\}$.
\lean{}\notready
\uses{lem:vertical-iff,def:fiber}
\end{lemma}

\begin{lemma}[Fibre composition is the group law]\label{lem:fiber-comp}
Composition in $p_{\Weier}.\mathtt{Fiber}\,(\Spec R)$ is the group law of
$\VC(R)$: the composite of $C:W\to W_0$ and $C':W_0\to W_1$ has $\VC$-component
$C'\cdot C$.
\lean{}\notready
\uses{lem:fiber-hom,def:total-comp,lem:vc-map-id}
\end{lemma}
\begin{proof}
By \eqref{eq:comp} with $f=g=\id_R$ the composite has $\VC$-component
$(\map\,\id\,C')\cdot C=C'\cdot C$, using Lemma~\ref{lem:vc-map-id}.
\end{proof}

\begin{lemma}[The fibre is a groupoid]\label{lem:fiber-groupoid}
$p_{\Weier}.\mathtt{Fiber}\,(\Spec R)$ is a groupoid: a morphism $C:W\to W_0$
(so $C\cdot W=W_0$) has two-sided inverse $C^{-1}:W_0\to W$, where
$C^{-1}\cdot W_0=W$.
\lean{}\notready
\uses{lem:fiber-hom,lem:fiber-comp,def:vc}
\end{lemma}
\begin{proof}
$C^{-1}\cdot W_0=C^{-1}\cdot(C\cdot W)=(C^{-1}C)\cdot W=W$ by \texttt{mul\_smul}
and \texttt{inv\_mul\_cancel}; the two composites are $C^{-1}C=\onel$ and
$CC^{-1}=\onel$ (Lemma~\ref{lem:fiber-comp}).
\end{proof}

\begin{lemma}[The fibre is the translation groupoid]\label{lem:fiber-descr}
For $R\in\Aff$ the fibre $p_{\Weier}.\mathtt{Fiber}\,(\Spec R)$ is isomorphic, as a
category, to the action groupoid $\VC(R)\ltimes\texttt{WeierstrassCurve}\,R$
(objects Weierstrass curves over $R$; $\Hom(W,W_0)=\{C\mid C\cdot W=W_0\}$;
composition the group law).  The isomorphism is the identity on objects and
$C\mapsto C$ on morphisms.
\lean{}\notready
\uses{lem:fiber-hom,lem:fiber-comp,lem:fiber-groupoid}
\end{lemma}

\begin{lemma}[Base-change data is a morphism]\label{lem:scart-welldef}
For $\underline f:\Spec R\to\Spec R'$ (ring map $f:R'\to R$) and
$(R',W')\in\Weier$, the pair $\lambda_{f,W'}\coloneqq(\underline f,\onel)$ is a
morphism $(R,\,W'.\map\,f)\to(R',W')$: indeed $\onel\cdot(W'.\map\,f)=W'.\map\,f$.
\lean{}\notready
\uses{def:total-hom,def:vc}
\end{lemma}

\begin{lemma}[Existence of the cartesian factorisation]\label{lem:scart-exist}
Let $(A,V)\in\Weier$, let $\underline h:\Spec A\to\Spec R$ (ring map $h:R\to A$),
and let $\psi=(\underline{k},D):(A,V)\to(R',W')$ lift
$\underline f\circ\underline h$, so $k=h\circ f$ and $D\cdot V=W'.\map\,k$.  Then
$\chi\coloneqq(\underline h,D)$ is a morphism $(A,V)\to(R,W'.\map\,f)$ and
$\chi\gg\lambda_{f,W'}=\psi$.
\lean{}\notready
\uses{lem:scart-welldef,def:total-comp,lem:map-func}
\end{lemma}
\begin{proof}
Legitimacy of $\chi$ is $D\cdot V=(W'.\map\,f).\map\,h$, which holds because
$D\cdot V=W'.\map\,(h\circ f)=(W'.\map\,f).\map\,h$ by Lemma~\ref{lem:map-func}.
By \eqref{eq:comp}, $\chi\gg\lambda_{f,W'}$ has $\Aff$-component
$\underline{h\circ f}=\underline k$ and $\VC$-component
$(\map\,h\,\onel)\cdot D=D$, so $\chi\gg\lambda_{f,W'}=\psi$.
\end{proof}

\begin{lemma}[Uniqueness of the cartesian factorisation]\label{lem:scart-uniq}
With notation as in Lemma~\ref{lem:scart-exist}, any morphism
$\chi'=(\underline h,E):(A,V)\to(R,W'.\map\,f)$ with
$\chi'\gg\lambda_{f,W'}=\psi$ has $E=D$; hence $\chi'=\chi$.
\lean{}\notready
\uses{def:total-comp}
\end{lemma}
\begin{proof}
By \eqref{eq:comp} the $\VC$-component of $\chi'\gg\lambda_{f,W'}$ is
$(\map\,h\,\onel)\cdot E=E$; equating with that of $\psi$ gives $E=D$.
\end{proof}

\begin{lemma}[Base-change arrows are strongly cartesian]\label{lem:scart}
For $\underline f:\Spec R\to\Spec R'$ and $(R',W')\in\Weier$, the arrow
$\lambda_{f,W'}=(\underline f,\onel):(R,\,W'.\map\,f)\to(R',W')$ is strongly
cartesian.
\lean{}\notready
\uses{def:cartesian,lem:scart-welldef,lem:scart-exist,lem:scart-uniq}
\end{lemma}
\begin{proof}
Lemmas~\ref{lem:scart-exist} and \ref{lem:scart-uniq} provide the unique
factorisation required by Definition~\ref{def:cartesian}(b).
\end{proof}

\begin{lemma}[$p_{\Weier}$ is prefibered]\label{lem:weier-prefibered}
$p_{\Weier}$ is prefibered: for every $\underline f:\Spec R\to\Spec R'$ and
$(R',W')$ over $\Spec R'$, the arrow $\lambda_{f,W'}$ (Lemma~\ref{lem:scart}) is a
strongly cartesian lift of $\underline f$ with target $(R',W')$.
\lean{}\notready
\uses{def:isfibered,lem:scart,def:proj}
\end{lemma}

\begin{theorem}[$p_{\Weier}$ is fibered]\label{thm:weier-fibered}
$p_{\Weier}:\Weier\to\Aff$ is a fibered category ($\isFib$).
\lean{}\notready
\uses{def:isfibered,lem:weier-prefibered,lem:scart-api}
\end{theorem}
\begin{proof}
$p_{\Weier}$ is prefibered (Lemma~\ref{lem:weier-prefibered}) and strongly
cartesian morphisms compose (Lemma~\ref{lem:scart-api}).
\end{proof}

\begin{proposition}[Groupoid fibres give a CFG]\label{prop:cfg-suff}
Let $p:\mathcal X\to\mathcal S$ be prefibered.  If every fibre
$p.\mathtt{Fiber}\,S$ is a groupoid, then every morphism of $\mathcal X$ is
strongly cartesian; i.e.\ $p$ is fibered in groupoids.
\lean{}\notready
\uses{def:cfg,def:fiber,def:isfibered}
\end{proposition}
\begin{proof}[Proof sketch]
Given $\varphi:a\to b$ over $f:R\to S$ and a test $\psi:a'\to b$ over $f\circ g$,
let $\lambda:f^*b\to b$ be a strongly cartesian lift (prefiberedness).  Then
$\psi=\chi_0\gg\lambda$ with $\chi_0$ over $g$, and $\varphi=v\gg\lambda$ with $v$
vertical, hence invertible as the fibre is a groupoid; $\chi\coloneqq\chi_0\gg
v^{-1}$ is the unique factorisation, so $\varphi$ is strongly cartesian.
\end{proof}

\begin{proposition}[A CFG has groupoid fibres]\label{prop:cfg-nec}
If $p:\mathcal X\to\mathcal S$ is fibered in groupoids, then every fibre is a
groupoid.
\lean{}\notready
\uses{def:cfg,def:fiber,lem:scart-api}
\end{proposition}
\begin{proof}[Proof sketch]
A fibre morphism is vertical, hence lifts the isomorphism $\onel_S$; being
strongly cartesian it is an isomorphism (Lemma~\ref{lem:scart-api}).
\end{proof}

\begin{proposition}[Groupoid criterion]\label{prop:cfg-criterion}
A prefibered $p:\mathcal X\to\mathcal S$ is fibered in groupoids iff every fibre
is a groupoid.
\lean{}\notready
\uses{prop:cfg-suff,prop:cfg-nec}
\end{proposition}

\begin{corollary}[$\Weier$ is a CFG]\label{cor:weier-cfg}
$p_{\Weier}$ is fibered in groupoids: it is prefibered
(Lemma~\ref{lem:weier-prefibered}) and each fibre is a groupoid
(Lemma~\ref{lem:fiber-groupoid}); apply Proposition~\ref{prop:cfg-suff}.
\lean{}\notready
\uses{lem:weier-prefibered,lem:fiber-groupoid,prop:cfg-suff}
\end{corollary}

\subsection{The pseudofunctor presentation (optional route)}

\begin{definition}[The base-change functor]\label{def:F-map}
For a ring map $f:R\to A$, the assignment $W\mapsto W.\map\,f$ on objects and
$C\mapsto\map\,f\,C$ on morphisms is a functor
$\mathbf F(f):\mathbf F(R)\to\mathbf F(A)$ between the fibre groupoids of
Lemma~\ref{lem:fiber-descr}, where
$\mathbf F(R)\coloneqq\VC(R)\ltimes\texttt{WeierstrassCurve}\,R$; functoriality on
morphisms is Lemma~\ref{lem:equivariance}.
\lean{}\notready
\uses{lem:fiber-descr,def:vc-map,lem:equivariance}
\end{definition}

\begin{construction}[The pseudofunctor $\mathbf F$]\label{con:pseudofunctor}
The data $R\mapsto\mathbf F(R)$, $f\mapsto\mathbf F(f)$
(Definition~\ref{def:F-map}) assemble into a pseudofunctor
$\mathbf F:\mathtt{LocallyDiscrete}\,\CRingCat\to\Cat$: the unitor
$\mathbf F(\id)\cong\mathrm{id}$ and the associator
$\mathbf F(g)\circ\mathbf F(f)\cong\mathbf F(g\circ f)$ are the identity $2$-cells
supplied by \texttt{map\_id} and \texttt{map\_map} (Lemma~\ref{lem:map-func}), and
the pentagon and triangle coherences hold because these $2$-cells are identities.
\lean{}\notready
\uses{def:F-map,lem:map-func}
\end{construction}

\begin{construction}[$\Weier$ as a CoGrothendieck construction]\label{con:weier-cogroth}
Under the identification $\Aff^{\op}=\CRingCat$, $\mathbf F$ is a contravariant
pseudofunctor on $\Aff$, and its CoGrothendieck construction
(Remark~\ref{rem:grothendieck}) satisfies $\int^{\mathrm c}\!\mathbf F\cong\Weier$
over $\Aff$, with the Mathlib projection
$\mathtt{forget}\,\mathbf F:\int^{\mathrm c}\!\mathbf F\to\Aff$ equal to
$p_{\Weier}$ (Definition~\ref{def:proj}) across this isomorphism.
\lean{CategoryTheory.Pseudofunctor.CoGrothendieck, CategoryTheory.Pseudofunctor.CoGrothendieck.forget}\notready
\uses{con:pseudofunctor,def:total,def:proj,rem:grothendieck}
\end{construction}
\begin{remark}
Along this route the fibered-category conclusions are supplied by Mathlib: by
Proposition~\ref{prop:cogroth-fibered}, $\mathtt{forget}\,\mathbf F$ carries an
\texttt{IsFibered} instance (with strongly cartesian lifts $\mathtt{cartesianLift}$),
which is Theorem~\ref{thm:weier-fibered} and the prefibered part of
Corollary~\ref{cor:weier-cfg} transported across
$\int^{\mathrm c}\!\mathbf F\cong\Weier$; since each fibre $\mathbf F(R)$ is a
groupoid (Lemma~\ref{lem:fiber-groupoid}), Proposition~\ref{prop:cfg-suff} gives
Corollary~\ref{cor:weier-cfg}.  The obligation this route incurs is
Construction~\ref{con:pseudofunctor}, i.e.\ producing the pseudofunctor
$\mathbf F$ with its coherences.  The direct route
(Definitions~\ref{def:total-obj}--Corollary~\ref{cor:weier-cfg}) instead proves
strong cartesianness by hand (Lemma~\ref{lem:scart}) and does not require the
bicategorical coherences.
\end{remark}

%==============================================================================
\section{The moduli stack of elliptic curves (Weierstrass model)}\label{sec:moduli}
%==============================================================================

\subsection{The moduli category and its fibered structure}

\begin{lemma}[$\isEll$ is stable under base change]\label{lem:iselliptic-map}
If $W.\isEll$ and $f:R\to A$ then $(W.\map\,f).\isEll$.  Indeed
$(W.\map\,f).\Delta=f(W.\Delta)$ (Lemma~\ref{lem:map-func}) is a unit, since a
ring homomorphism sends units to units.
\lean{WeierstrassCurve.IsElliptic.map}\notready
\uses{def:iselliptic,lem:map-func}
\end{lemma}

\begin{lemma}[$\isEll$ is stable under variable change]\label{lem:iselliptic-vc}
If $W.\isEll$ and $C=(u,r,s,t)\in\VC(R)$ then $(C\cdot W).\isEll$.  Indeed
$(C\cdot W).\Delta=u^{-12}\,W.\Delta$ (Lemma~\ref{lem:vc-invariants}) is a product
of units.
\lean{}\notready
\uses{def:iselliptic,lem:vc-invariants,def:vc}
\end{lemma}

\begin{definition}[The Weierstrass moduli category]\label{def:MW}
$\MW$ is the full subcategory of $\Weier$ on the objects $(R,W)$ with $W.\isEll$.
\lean{}\notready
\uses{def:total,def:iselliptic}
\end{definition}

\begin{definition}[The moduli projection]\label{def:MW-proj}
$p:\MW\to\Aff$ is the composite of the inclusion $\MW\hookrightarrow\Weier$ with
$p_{\Weier}$ (Definition~\ref{def:proj}).
\lean{}\notready
\uses{def:MW,def:proj}
\end{definition}

\begin{lemma}[Base-change lifts stay elliptic]\label{lem:MW-lift}
If $(R',W')\in\MW$ (so $W'.\isEll$) and $\underline f:\Spec R\to\Spec R'$, then the
source $(R,\,W'.\map\,f)$ of $\lambda_{f,W'}$ (Lemma~\ref{lem:scart}) lies in
$\MW$.
\lean{}\notready
\uses{def:MW,lem:scart,lem:iselliptic-map}
\end{lemma}

\begin{lemma}[The lifts are strongly cartesian in $\MW$]\label{lem:MW-scart}
For $(R',W')\in\MW$ and $\underline f:\Spec R\to\Spec R'$, the arrow
$\lambda_{f,W'}$ is strongly cartesian for $p:\MW\to\Aff$.
\lean{}\notready
\uses{lem:MW-lift,lem:scart,def:MW}
\end{lemma}
\begin{proof}
$\MW\subseteq\Weier$ is full, so its hom-sets between elliptic objects coincide
with those of $\Weier$.  The universal property (Definition~\ref{def:cartesian})
for $p$ quantifies over test objects of $\MW$, a subfamily of those for
$p_{\Weier}$; the unique factorisation produced in Lemma~\ref{lem:scart} has
source a test object of $\MW$ and target $(R,\,W'.\map\,f)\in\MW$
(Lemma~\ref{lem:MW-lift}), hence is a morphism of $\MW$.  Thus $\lambda_{f,W'}$ is
strongly cartesian in $\MW$.
\end{proof}

\begin{lemma}[$p:\MW\to\Aff$ is prefibered]\label{lem:MW-prefibered}
For every $\underline f:\Spec R\to\Spec R'$ and $(R',W')\in\MW$, the arrow
$\lambda_{f,W'}$ is a strongly cartesian lift of $\underline f$ in $\MW$ with
target $(R',W')$.
\lean{}\notready
\uses{def:isfibered,lem:MW-scart,def:MW-proj}
\end{lemma}

\begin{lemma}[The fibres of $\MW$ are groupoids]\label{lem:MW-fiber-groupoid}
For $R\in\Aff$, $p.\mathtt{Fiber}\,(\Spec R)$ is the full subcategory of
$p_{\Weier}.\mathtt{Fiber}\,(\Spec R)$ (Lemma~\ref{lem:fiber-descr}) on the
elliptic curves, and it is a groupoid.
\lean{}\notready
\uses{lem:fiber-groupoid,lem:iselliptic-vc,def:MW}
\end{lemma}
\begin{proof}
By Lemma~\ref{lem:iselliptic-vc}, if $W.\isEll$ and $C\cdot W=W_0$ then
$W_0.\isEll$; the fibre is thus the full subcategory of
Lemma~\ref{lem:fiber-descr} on elliptic objects, and the inverse $C^{-1}$
(Lemma~\ref{lem:fiber-groupoid}) preserves ellipticity, so it is a groupoid.
\end{proof}

\begin{theorem}[$\MW$ is a CFG]\label{thm:MW-cfg}
$p:\MW\to\Aff$ is fibered in groupoids.
\lean{}\notready
\uses{lem:MW-prefibered,lem:MW-fiber-groupoid,prop:cfg-suff}
\end{theorem}
\begin{proof}
$p$ is prefibered (Lemma~\ref{lem:MW-prefibered}) and its fibres are groupoids
(Lemma~\ref{lem:MW-fiber-groupoid}); apply Proposition~\ref{prop:cfg-suff}.
\end{proof}

\subsection{The \texorpdfstring{$j$}{j}-invariant}

\begin{definition}[The affine \texorpdfstring{$j$}{j}-line]\label{def:affine-line}
$\mathbb A^1$ is the CFG over $\Aff$ represented by $\mathbb Z[j]$: its fibre over
$R$ is the discrete category on the underlying set of $R$ (objects $(R,x)$ with
$x\in R$; only identity morphisms), and base change is application of the ring
map.
\lean{}\notready
\uses{def:cfg}
\end{definition}

\begin{lemma}[$j$ is invariant under variable change]\label{lem:j-vc}
For $W.\isEll$ and $C\in\VC(R)$, $(C\cdot W).j=W.j$.
\lean{WeierstrassCurve.variableChange\_j}\notready
\uses{def:iselliptic,lem:iselliptic-vc}
\end{lemma}

\begin{lemma}[$j$ is natural under base change]\label{lem:j-map}
For $W.\isEll$ and $f:R\to A$, $(W.\map\,f).j=f(W.j)$.
\lean{WeierstrassCurve.map\_j}\notready
\uses{def:iselliptic,lem:iselliptic-map}
\end{lemma}

\begin{definition}[The \texorpdfstring{$j$}{j}-line morphism]\label{def:jmap}
The assignment $(R,W)\mapsto(R,W.j)$ and $(\underline f,C)\mapsto\underline f$
defines a morphism of CFGs $\mathfrak j:\MW\to\mathbb A^1$ over $\Aff$:
Lemma~\ref{lem:j-vc} makes it well-defined on morphisms (the $\VC$-component acts
trivially on $j$), and Lemma~\ref{lem:j-map} makes it commute with base change.
\lean{}\notready
\uses{def:MW,def:affine-line,lem:j-vc,lem:j-map}
\end{definition}

\begin{lemma}[$\texttt{ofJ}$ realises a $j$-invariant]\label{lem:ofJ-j}
Over a field $F$, $(\texttt{ofJ}\,j_0).j=j_0$ for every $j_0\in F$.
\lean{WeierstrassCurve.ofJ\_j}\notready
\uses{def:iselliptic}
\end{lemma}

\begin{proposition}[Fibrewise essential surjectivity of $\mathfrak j$ over a field]\label{prop:j-surj}
Over a field $F$, $\mathfrak j$ is essentially surjective on the fibre: every
$j_0\in F$ equals $W.j$ for $W=\texttt{ofJ}\,j_0$ (Lemma~\ref{lem:ofJ-j}).  The
complementary fibrewise injectivity holds over separably closed fields and is the
target \texttt{isom\_iff\_j\_eq}.
\lean{}\notready
\uses{def:jmap,lem:ofJ-j}
\end{proposition}

\subsection{The quotient presentation}

\begin{definition}[Universal Weierstrass ring]\label{def:univ-ring-P}
$P\coloneqq\mathbb Z[a_1,a_2,a_3,a_4,a_6]=\texttt{MvPolynomial (Fin 5)}\ \mathbb Z$,
the polynomial ring on the five Weierstrass coefficients.
\lean{}\notready
\uses{def:weier}
\end{definition}

\begin{definition}[Tautological curve and its discriminant]\label{def:univ-curve-P}
$W^{P}:\texttt{WeierstrassCurve}\,P$ is the curve whose coefficients are the five
generators of $P$; set $\Delta_P\coloneqq W^{P}.\Delta\in P$.
\lean{}\notready
\uses{def:univ-ring-P,def:weier}
\end{definition}

\begin{definition}[Universal base ring]\label{def:univ-ring-A}
$A\coloneqq P[\Delta_P^{-1}]=\texttt{Localization.Away}\,\Delta_P$, with structure
map $\iota:P\to A$.
\lean{}\notready
\uses{def:univ-curve-P}
\end{definition}

\begin{definition}[Universal elliptic curve]\label{def:univ}
$W^{\mathrm{univ}}\coloneqq W^{P}.\map\,\iota:\texttt{WeierstrassCurve}\,A$.
\lean{}\notready
\uses{def:univ-ring-A,def:univ-curve-P,def:weier}
\end{definition}

\begin{lemma}[The universal curve is elliptic]\label{lem:univ-elliptic}
$W^{\mathrm{univ}}.\isEll$, because
$\iota(\Delta_P)=W^{\mathrm{univ}}.\Delta$ is a unit in
$A=\texttt{Localization.Away}\,\Delta_P$.
\lean{}\notready
\uses{def:univ,def:iselliptic,lem:iselliptic-map}
\end{lemma}

\begin{definition}[The group functor of variable changes]\label{def:G}
$G\colon\CRing\to\mathrm{Grp}$ is the functor $R\mapsto\VC(R)$, with the group law
of Definition~\ref{def:vc} and functoriality $\map$ (Definition~\ref{def:vc-map}).
\lean{}\notready
\uses{def:vc,def:vc-map}
\end{definition}

\begin{lemma}[Representability of $G$]\label{lem:G-representable}
$G$ is representable by $\mathbb Z[u,u^{-1},r,s,t]$; as a scheme
$G\cong\Gm\times\Ga^3$, with group law the semidirect product $\Ga^3\rtimes\Gm$
for the evident weights.
\lean{}\notready
\uses{def:G,def:vc}
\end{lemma}

\begin{definition}[The $G$-action on $\Spec A$]\label{def:G-action}
$G$ acts on $\Spec A$ by variable change, $C\cdot W^{\mathrm{univ}}$; the co-action
is the $A$-algebra map $A\to A\otimes\mathbb Z[u,u^{-1},r,s,t]$ dual to the
transformed coefficients of Definition~\ref{def:vc}.
\lean{}\notready
\uses{def:G,def:univ,def:vc}
\end{definition}

\begin{lemma}[Points of $\Spec A$ are elliptic curves]\label{lem:Apoint-elliptic}
For $R\in\CRing$, the assignment $(\varphi:A\to R)\mapsto W^{\mathrm{univ}}.\map\,\varphi$
is a bijection from $\Hom_{\CRing}(A,R)$ to the set of elliptic
$W:\texttt{WeierstrassCurve}\,R$.
\lean{}\notready
\uses{def:univ,def:univ-ring-A,def:iselliptic}
\end{lemma}
\begin{proof}
By the universal property of $P=\texttt{MvPolynomial (Fin 5)}\,\mathbb Z$, a ring
map $P\to R$ is a choice of $(a_i)\in R^5$; by the universal property of
$\texttt{Localization.Away}\,\Delta_P$ it extends to $A\to R$ iff the image of
$\Delta_P$, namely $W.\Delta$, is a unit, i.e.\ iff $W.\isEll$
(Definition~\ref{def:iselliptic}).
\end{proof}

\begin{lemma}[The bijection is $G$-equivariant and natural]\label{lem:quotient-compat}
Under Lemma~\ref{lem:Apoint-elliptic}, the $G(R)=\VC(R)$-action on
$\Hom_{\CRing}(A,R)$ (Definition~\ref{def:G-action}) corresponds to the
variable-change action on elliptic curves, and the bijection is natural in $R$
along base change (Lemma~\ref{lem:equivariance}).
\lean{}\notready
\uses{lem:Apoint-elliptic,def:G-action,lem:equivariance}
\end{lemma}

\begin{theorem}[Quotient presentation]\label{thm:quotient}
The universal curve induces an equivalence of CFGs over $\Aff$
\[
   \MW \;\simeq\; [\,\Spec A\,/\,G\,],
\]
where $[\Spec A/G]$ is the quotient CFG whose fibre over $R$ is the action
groupoid $G(R)\ltimes(\Spec A)(R)$.
\lean{}\notready
\uses{thm:MW-cfg,lem:univ-elliptic,lem:G-representable,lem:Apoint-elliptic,lem:quotient-compat,lem:fiber-descr}
\end{theorem}
\begin{proof}[Proof blueprint]
On each fibre, Lemma~\ref{lem:Apoint-elliptic} identifies the objects
$(\Spec A)(R)$ of $[\Spec A/G]$ with the elliptic curves over $R$, and
Lemma~\ref{lem:quotient-compat} identifies the $G(R)$-action with the
variable-change action; hence the action groupoid $G(R)\ltimes(\Spec A)(R)$ agrees
with the fibre groupoid of $\MW$ (Lemma~\ref{lem:fiber-descr} on elliptic curves).
Naturality in $R$ (Lemma~\ref{lem:quotient-compat}) upgrades these fibre
equivalences to an equivalence of CFGs.
\end{proof}

\begin{remark}
Theorem~\ref{thm:quotient} exhibits $\MW$ as a global quotient stack.  It is the
entry point to the smoothness of $\Mss$ (the atlas $\Spec A$ is smooth over
$\mathbb Z$) and to the automorphism group schemes (stabilisers of the
$G$-action), including the generic automorphism $[-1]$ and the extra automorphisms
at $j=0,1728$.
\end{remark}

%==============================================================================
\section{Descent and the stack condition}\label{sec:descent}
%==============================================================================

We now formulate what it means for the CFG $\MW$ to be a stack for
$\tau_{\mathrm{fppf}}$, and reduce it to descent statements.  As flagged in
Warning~\ref{warn:no-stacks}, the notion of stack is not in Mathlib; the
definitions here are new.

\begin{definition}[Isom presheaf]\label{def:isompre}
For a CFG $p:\mathcal X\to\mathcal S$, $S\in\mathcal S$ and $x,y\in\mathcal X(S)$,
the presheaf $\Isom_S(x,y)$ on the slice $\mathcal S_{/S}$ sends
$(T\xrightarrow{a}S)$ to the set of isomorphisms $a^*x\xrightarrow{\sim}a^*y$ in
$\mathcal X(T)$, with restriction by further base change.
\lean{}\notready
\uses{def:cfg}
\end{definition}

\begin{definition}[Prestack, stack]\label{def:stack}
Let $(\mathcal S,\tau)$ be a site and $p:\mathcal X\to\mathcal S$ a CFG.
\begin{enumerate}[label=(\alph*)]
\item $p$ is a \emph{prestack} for $\tau$ if for all $S$ and all
$x,y\in\mathcal X(S)$ the presheaf $\Isom_S(x,y)$ is a $\tau$-sheaf on
$\mathcal S_{/S}$ (equivalently: morphisms glue).
\item A \emph{descent datum} for a covering $\{U_i\to S\}$ is a family
$x_i\in\mathcal X(U_i)$ with isomorphisms
$\varphi_{ij}:x_i|_{U_{ij}}\xrightarrow{\sim}x_j|_{U_{ij}}$ satisfying the cocycle
condition on triple overlaps.  It is \emph{effective} if it is isomorphic to the
descent datum induced by some $x\in\mathcal X(S)$.
\item $p$ is a \emph{stack} for $\tau$ if it is a prestack and every descent
datum for every $\tau$-covering is effective.
\end{enumerate}
\lean{}\notready
\uses{def:isompre,def:site}
\end{definition}

\begin{proposition}[Representability of $\Isom$ for $\MW$]\label{prop:isom-rep}
For $W_1,W_2$ elliptic over $R$, the presheaf $\Isom_{\Spec R}(W_1,W_2)$ is
representable by an affine $R$-scheme: a closed subscheme of
$G_R=\Gm\times\Ga^3$ (base-changed to $R$) cut out by the five equations
$C\cdot W_1=W_2$ of Definition~\ref{def:vc} (coefficient-wise).  In particular it
is an fppf sheaf.
\lean{}\notready
\uses{def:isompre,def:vc,def:MW}
\end{proposition}
\begin{proof}[Proof blueprint]
A $T$-point of $\Isom(W_1,W_2)$ over $a:R\to T$ is, by
Lemma~\ref{lem:fiber-descr}, a $C\in\VC(T)$ with
$C\cdot(W_1.\map\,a)=W_2.\map\,a$.  Equivalently it is a $T$-point of
$G_R=\Spec R[u,u^{-1},r,s,t]$ satisfying the five polynomial equations obtained by
equating coefficients in $C\cdot W_1=W_2$ (a system with $\mathbb Z$-coefficients
in $u^{\pm1},r,s,t$ and the $a_i(W_1),a_i(W_2)$).  Hence
$\Isom(W_1,W_2)=\Spec\bigl(R[u,u^{-1},r,s,t]/(\text{those }5\text{ relations})
\bigr)$, an affine scheme over $R$; representable functors on $\Aff$ are fppf
sheaves once faithfully flat descent for morphisms is available (see
Lemma~\ref{lem:repr-sheaf}).
\end{proof}

\begin{lemma}[Representable $\Rightarrow$ fppf sheaf]\label{lem:repr-sheaf}
The functor of points of any affine scheme is a sheaf for $\tau_{\mathrm{fppf}}$
on $\Aff$.  Equivalently: for a faithfully flat ring map $R\to R'$ the Amitsur
complex
\[
   0\to R\to R'\rightrightarrows R'\otimes_R R'\to\cdots
\]
is exact; hence $R\to R'$ is an equaliser of
$R'\rightrightarrows R'\otimes_R R'$ and any $R$-algebra of finite type has the
descent property.
\lean{RingHom.FaithfullyFlat, Algebra.amitsurComplex (target)}\notready
\uses{def:fppf-cover,def:site}
\end{lemma}
\begin{proof}[Status]
The exactness of the Amitsur (Čech--Amitsur) complex for a faithfully flat
$R\to R'$ is the key input.  Mathlib has \texttt{Module.FaithfullyFlat} and the
one-step statement that $R\to R'$ is injective with the expected equaliser
description in low degree; the full cosimplicial exactness and the packaging
``representable $\Rightarrow$ fppf sheaf'' are targets.  This lemma is a
self-contained prerequisite, independent of elliptic curves.
\end{proof}

\begin{theorem}[$\MW$ is an fppf prestack]\label{thm:prestack}
$p:\MW\to\Aff$ is a prestack for $\tau_{\mathrm{fppf}}$.
\lean{}\notready
\uses{def:stack,prop:isom-rep,lem:repr-sheaf}
\end{theorem}
\begin{proof}
Immediate from Proposition~\ref{prop:isom-rep} and Lemma~\ref{lem:repr-sheaf}:
each $\Isom$ presheaf is representable, hence an fppf sheaf, which is precisely
the prestack condition of Definition~\ref{def:stack}(a).
\end{proof}

\begin{theorem}[Effectivity of descent for $\MW$]\label{thm:effective}
Every fppf descent datum for $\MW$ is effective; hence
$p:\MW\to\Aff$ is a stack for $\tau_{\mathrm{fppf}}$.
\lean{}\notready
\uses{def:stack,thm:prestack}
\end{theorem}
\begin{proof}[Proof blueprint and prerequisite ledger]
Let $\{R\to R_i\}$ be an fppf cover with $R\to R'=\prod R_i$ faithfully flat.  A
descent datum is a curve $W'$ over $R'$ together with an isomorphism
$\sigma:\mathrm{pr}_1^*W'\xrightarrow{\sim}\mathrm{pr}_2^*W'$ over
$R'\otimes_R R'$ satisfying the cocycle condition over
$R'\otimes_R R'\otimes_R R'$.  Effectivity means: there is an elliptic
$W$ over $R$ and an isomorphism $W'\cong W.\map\,(R\to R')$ carrying the canonical
datum to $(W',\sigma)$.  The reduction proceeds by descending the coefficients:
\begin{enumerate}[label=(\roman*)]
\item Each coefficient $a_k(W')\in R'$ together with the transformation
prescribed by $\sigma$ (a variable change $C\in\VC(R'\otimes_R R')$) constitutes
a descent datum for the quasi-coherent sheaf underlying the Weierstrass data.
Because $\sigma$ need not be the identity, the naive coefficient-wise descent
fails in general; what descends is the pair (line bundle $\omega$, section data),
which is why the twist by $\mathrm{Pic}$ appears (Warning~\ref{warn:pic}).
\item \textbf{Prerequisite: fppf descent for quasi-coherent modules.}  For a
faithfully flat $R\to R'$, the category of $R$-modules is equivalent to the
category of $R'$-modules with descent data (Grothendieck).  This is
\emph{not} in Mathlib; it is the load-bearing missing input.  Its statement:
the base-change functor $\mathrm{Mod}_R\to\mathrm{Desc}(R'/R)$ is an equivalence.
\item \textbf{Prerequisite: fppf descent for the multiplicative and additive
group.}  $\Gm$ and $\Ga$ are fppf sheaves (special case of (ii) / of
Lemma~\ref{lem:repr-sheaf}); this handles the descent of the unit $u$ and of the
translation parameters $r,s,t$ in $\sigma$.
\item Granting (ii)--(iii), the descent datum $(W',\sigma)$ is effective: the
Hodge line bundle $\omega'=p'_*\Omega$ (in the affine model, the invertible
$R'$-module attached to $\sigma$) descends to an invertible $R$-module $\omega$;
choosing (fppf-locally) a trivialisation recovers a Weierstrass equation, and the
cocycle condition guarantees the descended $(a_k)\in R$ define an elliptic curve
$W$ over $R$ inducing $(W',\sigma)$.  When $\omega$ is globally trivial (e.g.\
$\mathrm{Pic}(R)=0$) the descent is a genuine Weierstrass curve over $R$; in
general $W$ exists in the stackification.
\end{enumerate}
Thus, modulo the descent prerequisites (ii)--(iii), $\MW$ satisfies effective
descent and is an fppf stack.  The honest statement of current feasibility: the
prestack theorem~\ref{thm:prestack} is reachable from present Mathlib plus
Lemma~\ref{lem:repr-sheaf}; the effectivity theorem requires the quasi-coherent
descent package (ii), a substantial independent development.
\end{proof}

\begin{remark}[Stackification]\label{rem:stackification}
Since Mathlib lacks stackification, the clean statement ``the fppf-stackification
of $\MW$ is $\Mss$'' (Theorem~\ref{thm:comparison}) presupposes a new
\texttt{Stackification} construction (the $2$-categorical analogue of
\texttt{presheafToSheaf}).  Building it is a self-contained category-theory
target: sheafify the object- and morphism-presheaves and check the universal
property.  It is listed as milestone \textbf{M8}.
\end{remark}

%==============================================================================
\section{Comparison with the geometric moduli stack}\label{sec:comparison}
%==============================================================================

\begin{definition}[Elliptic curve over a scheme; $\Mss$]\label{def:geomell}
An \emph{elliptic curve} over a scheme $S$ is a morphism $p:E\to S$ that is proper,
smooth, of relative dimension one, with geometrically connected fibres of
arithmetic genus one, together with a section $e:S\to E$.  A morphism of elliptic
curves over $a:T\to S$ is a cartesian square inducing $E_T\xrightarrow{\sim}
a^*E$ compatibly with sections.  The resulting category $\Mss$ over $\Sch$ (or
over $\Aff$, or over $\Sch_{/\mathbb Z}$) is fibered in groupoids; this is the
\emph{moduli stack of elliptic curves} \cite{DeligneRapoport,KatzMazur,Olsson}.
\lean{}\notready
\uses{def:cfg}
\end{definition}

\begin{theorem}[Comparison]\label{thm:comparison}
There is a morphism of CFGs $\Phi:\MW\to\Mss$ over $\Aff$ sending an elliptic
Weierstrass curve $W$ over $R$ to $\mathrm{Proj}$ of its homogeneous coordinate
ring (the projective plane cubic with the point at infinity as section).  Then:
\begin{enumerate}[label=(\alph*)]
\item For $R$ with $\mathrm{Pic}(R)[12]=0$ (in particular $\mathrm{Pic}(R)=0$,
e.g.\ $R$ local or a PID), $\Phi_R:\MW(R)\to\Mss(R)$ is an equivalence of
groupoids.
\item In general $\Phi$ exhibits $\Mss$ as the fppf-stackification of $\MW$;
equivalently, every elliptic curve over $S$ admits a Weierstrass equation
Zariski-locally on $S$, and the obstruction to a global one lies in
$\mathrm{Pic}(S)$ via the Hodge bundle.
\end{enumerate}
\lean{}\notready
\uses{def:MW,def:geomell,thm:quotient,rem:stackification,warn:pic}
\end{theorem}
\begin{proof}[Discussion and prerequisites]
The construction of $\Phi$ and the proof of (a),(b) are the content of
\cite[II.2, III.1, III.3]{Silverman}, \cite[\S2.1--2.2]{Deligne1975},
\cite[Ch.~2]{KatzMazur}.  The essential geometric facts are: (1) an elliptic
curve $E/S$ has $p_*\mathcal O_E=\mathcal O_S$ and an invertible
$\omega=p_*\Omega_{E/S}$; (2) Riemann--Roch for the section divisor yields a
closed immersion $E\hookrightarrow\mathbb P(p_*\mathcal O_E(3e))$ realising $E$
as a Weierstrass cubic once $\omega$ is trivialised; (3) two Weierstrass models of
the same $E$ differ by a variable change, matching
Definition~\ref{def:vc}.  Each of these presupposes machinery not in Mathlib.

\medskip\noindent\textbf{Missing scheme-theoretic prerequisites
(future-work ledger).}  In rough dependency order:
\begin{enumerate}[label=(P\arabic*),leftmargin=3.2em]
\item proper morphisms of schemes and their stability
      (\texttt{AlgebraicGeometry.IsProper}, partial);
\item smooth morphisms and relative dimension (partial in
      \texttt{AlgebraicGeometry.Morphisms.Smooth});
\item relative differentials $\Omega_{E/S}$ and the dualising/Hodge line bundle;
\item higher direct images $R^i p_*$ and cohomology-and-base-change, giving
      $p_*\mathcal O_E=\mathcal O_S$ and invertibility of $\omega$;
\item relative Riemann--Roch / relative $\mathrm{Proj}$ and the very-ampleness of
      $\mathcal O_E(3e)$;
\item the classification of Weierstrass models of a fixed $E$ (the ``$3e$''
      embedding and its normal form), reproducing Definition~\ref{def:vc};
\item faithfully flat / fppf descent for quasi-coherent sheaves
      (Theorem~\ref{thm:effective}(ii)) and for schemes;
\item the $2$-categorical stackification of Remark~\ref{rem:stackification}.
\end{enumerate}
Consequently, within a horizon fixed by present Mathlib, the achievable target is
Parts~\ref{sec:weier}--\ref{sec:moduli} (the CFG $\MW$, its groupoid fibres, and
the quotient presentation) together with Theorem~\ref{thm:prestack} (the fppf
prestack property, modulo Lemma~\ref{lem:repr-sheaf}).  Theorem~\ref{thm:comparison}
and Theorem~\ref{thm:effective} are longer-term and are stated here to fix the
mathematically correct target and to make the dependency on (P1)--(P8) explicit.
\end{proof}

\begin{remark}[Algebraicity of $\Mss$]\label{rem:algebraicity}
With (P1)--(P8) in hand, Theorem~\ref{thm:quotient} shows $\Mss$ is a global
quotient stack $[\Spec A/G]$, hence a smooth algebraic (Artin) stack of relative
dimension one over $\Spec\mathbb Z$, separated, and Deligne--Mumford over
$\mathbb Z[1/6]$ (the automorphism group schemes are finite étale away from
characteristics $2,3$) \cite{DeligneRapoport,Olsson}.  These are statements about
$\Mss$, downstream of the entire ledger, recorded for orientation.
\end{remark}

%==============================================================================
\section{Milestone ladder}\label{sec:ladder}
%==============================================================================

The following orders the work by Mathlib-readiness.  Milestones
\textbf{M0}--\textbf{M5} are reachable from present Mathlib; \textbf{M6}--\textbf{M9}
depend on independent developments (descent, stackification, scheme geometry).

\begin{description}[leftmargin=3.4em,style=nextline]
\item[\textbf{M0} — Substrate audit.]  Confirm the interface of
  Definitions~\ref{def:homlift}--\ref{def:site}, \ref{def:weier}--\ref{def:iselliptic}
  against the pinned Mathlib commit; add missing thin API
  (\texttt{VariableChange.map}, \texttt{map\_j}, \texttt{IsElliptic.map} if
  absent).  \emph{Depends on:} nothing.
\item[\textbf{M1} — Equivariance.]  Definition~\ref{def:vc-map};
  Lemmas~\ref{lem:vc-map-one}, \ref{lem:vc-map-mul}, \ref{lem:vc-map-id},
  \ref{lem:vc-map-comp}; Lemma~\ref{lem:equivariance}.  \emph{Depends on:} M0.
\item[\textbf{M2} — The fibered category $\Weier$.]
  Definitions~\ref{def:total-obj}--\ref{def:total-comp} with
  Lemmas~\ref{lem:id-welldef}, \ref{lem:comp-welldef}; the category axioms
  \ref{lem:total-id-comp}--\ref{lem:total-assoc} and Definition~\ref{def:total};
  the projection \ref{def:proj}; the lift/fibre lemmas
  \ref{lem:homlift-mor}--\ref{lem:fiber-descr}; strong cartesianness
  \ref{lem:scart-welldef}--\ref{lem:scart}; prefiberedness
  \ref{lem:weier-prefibered}; Theorem~\ref{thm:weier-fibered}; the groupoid
  criterion \ref{prop:cfg-suff}--\ref{prop:cfg-criterion}; Corollary~\ref{cor:weier-cfg}.
  \emph{Depends on:} M1 and the \texttt{FiberedCategory} API.
\item[\textbf{M3} — The moduli CFG $\MW$.]  Lemmas~\ref{lem:iselliptic-map},
  \ref{lem:iselliptic-vc}; Definitions~\ref{def:MW}, \ref{def:MW-proj};
  Lemmas~\ref{lem:MW-lift}--\ref{lem:MW-fiber-groupoid}; Theorem~\ref{thm:MW-cfg};
  the $j$-invariant nodes \ref{def:affine-line}, \ref{lem:j-vc}, \ref{lem:j-map},
  \ref{def:jmap}, \ref{lem:ofJ-j}, \ref{prop:j-surj}.  \emph{Depends on:} M2.
\item[\textbf{M4} — Quotient presentation.]
  Definitions~\ref{def:univ-ring-P}--\ref{def:univ} with
  Lemma~\ref{lem:univ-elliptic}; the group nodes \ref{def:G},
  \ref{lem:G-representable}, \ref{def:G-action}; Lemmas~\ref{lem:Apoint-elliptic},
  \ref{lem:quotient-compat}; Theorem~\ref{thm:quotient}.  \emph{Depends on:} M3
  and the functor-of-points/Yoneda API for \texttt{MvPolynomial} and
  \texttt{Localization.Away}.
\item[\textbf{M5} — Pseudofunctor route (optional).]
  Definition~\ref{def:F-map}; Constructions~\ref{con:pseudofunctor},
  \ref{con:weier-cogroth}; along it Theorem~\ref{thm:weier-fibered} is discharged
  by Proposition~\ref{prop:cogroth-fibered}.  \emph{Depends on:} M1 (to build the
  pseudofunctor $\mathbf F$, including its coherences) and
  \texttt{Pseudofunctor.CoGrothendieck} together with
  \texttt{FiberedCategory.Grothendieck}
  (\texttt{CoGrothendieck.instIsFiberedForget}).
\item[\textbf{M6} — fppf site.]  Definition~\ref{def:fppf-cover};
  Lemma~\ref{lem:fppf-pretop}; Lemma~\ref{lem:repr-sheaf}.  \emph{Depends on:}
  flatness/finite-presentation base-change API and the Amitsur complex.
\item[\textbf{M7} — Prestack.]  Definitions~\ref{def:isompre},\ref{def:stack};
  Proposition~\ref{prop:isom-rep}; Theorem~\ref{thm:prestack}.  \emph{Depends on:}
  M3, M6.
\item[\textbf{M8} — Descent and stackification.]  Theorem~\ref{thm:effective}
  (needs quasi-coherent fppf descent, P7); Remark~\ref{rem:stackification} (the
  stackification functor).  \emph{Depends on:} M7 and P7.
\item[\textbf{M9} — Geometric comparison.]  Definition~\ref{def:geomell};
  Theorem~\ref{thm:comparison}; Remark~\ref{rem:algebraicity}.  \emph{Depends on:}
  M8 and the scheme-geometry ledger (P1)--(P6).
\end{description}

%==============================================================================
\section*{Appendix A. Explicit formulas and their Mathlib names}
\addcontentsline{toc}{section}{Appendix A. Explicit formulas}
%==============================================================================

The variable-change action is that of Definition~\ref{def:vc}, matching
\cite[III, Table~3.1]{Silverman} and the Mathlib instance
\texttt{WeierstrassCurve.instMulActionVariableChange}.  The transformation weights
are: $c_4\mapsto u^{-4}c_4$, $c_6\mapsto u^{-6}c_6$, $\Delta\mapsto
u^{-12}\Delta$, and $j\mapsto j$ (invariant); the corresponding Mathlib lemmas are
\texttt{variableChange\_c4}, \texttt{variableChange\_c6},
\texttt{variableChange\_Delta}, \texttt{variableChange\_j}.  Under base change
$f:R\to A$ the invariants are natural: $c_4,c_6,\Delta,j$ commute with $f$, via
\texttt{map\_c4}, \texttt{map\_c6}, \texttt{map\_Delta}, \texttt{map\_j}.  The
identity $c_4^3-c_6^2=1728\,\Delta$ is \texttt{WeierstrassCurve.c4\_cubed\_...}
(coefficient identity in \texttt{Mathlib.AlgebraicGeometry.EllipticCurve.Weierstrass}).

%==============================================================================
\section*{Appendix B. Formalisation status by item}
\addcontentsline{toc}{section}{Appendix B. Formalisation status}
%==============================================================================

\noindent\textbf{Already in Mathlib} (mark \leanok): the fibered-category API
(Definitions~\ref{def:homlift}--\ref{def:fiber},
Lemma~\ref{lem:scart-api}, Definition~\ref{def:isfibered}); the Grothendieck and
CoGrothendieck constructions of a pseudofunctor (Remark~\ref{rem:grothendieck})
and the theorem that the CoGrothendieck projection is a fibered category
(Proposition~\ref{prop:cogroth-fibered},
\texttt{Mathlib.CategoryTheory.FiberedCategory.Grothendieck}); sites and
sheaves (Definition~\ref{def:site}); the Weierstrass/elliptic-curve API
(Definitions~\ref{def:weier}--\ref{def:iselliptic},
Lemmas~\ref{lem:map-func},\ref{lem:vc-invariants}).

\smallskip\noindent\textbf{New but reachable now} (M1--M5, M6-partial):
Lemma~\ref{lem:equivariance}; the category $\Weier$ and its CFG structure;
$\MW$ and Theorem~\ref{thm:MW-cfg}; the quotient presentation
Theorem~\ref{thm:quotient}; the fppf pretopology and
Lemma~\ref{lem:repr-sheaf} (the latter needing the Amitsur complex).

\smallskip\noindent\textbf{New and dependent on major developments} (M8--M9):
Theorem~\ref{thm:effective} (quasi-coherent fppf descent, P7); stackification
(Remark~\ref{rem:stackification}); Definition~\ref{def:geomell} and
Theorem~\ref{thm:comparison} (the scheme-geometry ledger P1--P8);
Remark~\ref{rem:algebraicity}.

\begin{thebibliography}{9}

\bibitem{Deligne1975}
P.~Deligne, \emph{Courbes elliptiques: formulaire d'après J.~Tate}, in
\emph{Modular Functions of One Variable IV}, Lecture Notes in Math.\ \textbf{476},
Springer, 1975, 53--73.

\bibitem{DeligneRapoport}
P.~Deligne, M.~Rapoport, \emph{Les schémas de modules de courbes elliptiques}, in
\emph{Modular Functions of One Variable II}, Lecture Notes in Math.\ \textbf{349},
Springer, 1973, 143--316.

\bibitem{FGA}
B.~Fantechi et al., \emph{Fundamental Algebraic Geometry: Grothendieck's FGA
Explained}, Math.\ Surveys Monogr.\ \textbf{123}, AMS, 2005.  (Vistoli's chapter
on Grothendieck topologies, fibered categories and descent.)

\bibitem{KatzMazur}
N.~Katz, B.~Mazur, \emph{Arithmetic Moduli of Elliptic Curves}, Ann.\ of Math.\
Stud.\ \textbf{108}, Princeton, 1985.

\bibitem{Olsson}
M.~Olsson, \emph{Algebraic Spaces and Stacks}, AMS Colloq.\ Publ.\ \textbf{62},
2016.

\bibitem{Silverman}
J.~H.~Silverman, \emph{The Arithmetic of Elliptic Curves}, 2nd ed., GTM
\textbf{106}, Springer, 2009.  (Chapter III, especially Table~3.1.)

\bibitem{Vistoli2008}
A.~Vistoli, \emph{Notes on Grothendieck topologies, fibered categories and descent
theory}, 2008.

\bibitem{StacksProject}
The Stacks Project Authors, \emph{The Stacks Project},
\texttt{https://stacks.math.columbia.edu}.

\end{thebibliography}

\end{document}
